% Generated draft: requires mathematical and source review.
\documentclass[11pt]{article}
\usepackage[margin=25mm]{geometry}
\usepackage{fontspec}
\setmainfont{Latin Modern Roman}
\usepackage{amsmath,amssymb,mathtools,unicode-math}
\setmathfont{Latin Modern Math}
\usepackage{enumitem,needspace,xurl,hyperref}
\setlength{\parindent}{0pt}
\setlength{\parskip}{5pt}
\setlength{\emergencystretch}{3em}
\newcommand{\N}{\mathbb{N}}
\newcommand{\Z}{\mathbb{Z}}
\newcommand{\Q}{\mathbb{Q}}
\newcommand{\R}{\mathbb{R}}
\newcommand{\C}{\mathbb{C}}
\newcommand{\F}{\mathbb{F}}
\newcommand{\A}{\mathbb{A}}
\newcommand{\PP}{\mathbb{P}}
\newcommand{\E}{\mathbb{E}}
\newcommand{\eps}{\varepsilon}
\newcommand{\dd}{\,\mathrm{d}}
\begin{document}
\section*{UnsolvedMath ID: 701149}
\textbf{Source problem code:} catalogue entry
\subsection{Definitions and mathematical statement}
The supplied question is recorded below under the usual finite-group interpretation. The source excerpt does not explicitly specify finiteness, the permitted sets of primes, or conventions for trivial groups; these qualifications require review.

Let \(\pi\) be a nonempty set of primes, and let \(\pi'\) be its complement in the set of all primes. A finite \(\pi\)-group is a finite group whose order has prime divisors only in \(\pi\). Interpret a finite group \(G\) as \(\pi\)-soluble if it has a normal series whose factors are either \(\pi'\)-groups or \(p\)-groups for primes \(p\in\pi\). A Hall \(\pi\)-subgroup \(G_\pi\leq G\) satisfies that \(|G_\pi|\) is a \(\pi\)-number and \([G:G_\pi]\) is a \(\pi'\)-number.

For a finite soluble group \(H\), let \(n(H)\) denote its nilpotent length: the least number of nontrivial factors in a normal series of \(H\) all of whose factors are nilpotent. For a finite \(p\)-soluble group \(G\), let \(l_p(G)\) denote its ordinary \(p\)-length, the least number of nontrivial \(p\)-group factors in a normal series whose factors are either \(p\)-groups or \(p'\)-groups. Define the nilpotent \(\pi\)-length \(l_\pi^n(G)\) to be the least number of nontrivial nilpotent \(\pi\)-group factors in a normal series of \(G\) whose factors are either \(\pi'\)-groups or nilpotent \(\pi\)-groups. The superscript \(n\) is part of the notation for this invariant, not an exponent.

\textbf{Question (V. S. Monakhov; supplied source formulation).} Is it true that every \(\pi\)-soluble group \(G\), with a Hall \(\pi\)-subgroup \(G_\pi\), satisfies
\[
 l_\pi^n(G)\leq n(G_\pi)-1+\max_{p\in\pi}l_p(G)?
\]
The finite-group definitions above give a conventional reading of this question, subject to the unresolved boundary qualifications below.

The supplied source reports an affirmative answer when every proper subgroup of \(G_\pi\) is supersoluble. Here a finite group is supersoluble if it admits a normal series with cyclic factors of prime order. This affirmative special case is a supplied source claim, not an independently verified result.

\textbf{Unresolved boundary qualification.} With the usual empty-series conventions, \(n(1)=0\) and the length invariants of the trivial group are zero. The source does not explain how its displayed inequality is intended to apply when \(G_\pi=1\), nor how to interpret the maximum if \(\pi=\varnothing\). No additional hypothesis or exceptional convention resolving these cases is supplied.
\subsection{Short English statement}
Monakhov asks whether the nilpotent \(\pi\)-length of a \(\pi\)-soluble group is at most the nilpotent length of a Hall \(\pi\)-subgroup, minus one, plus the largest ordinary \(p\)-length for \(p\in\pi\). The source reports an affirmative answer when every proper subgroup of that Hall subgroup is supersoluble. Finiteness and degenerate cases require clarification.
\subsection{Sources}
\begin{itemize}
\item Local UnsolvedMath registry, permanent ID \(701149\), exact source problem code \texttt{catalogue entry}, title \emph{Kourovka Notebook Problem 15.61}; supplied complete source record, upstream revision \texttt{d371ee917ac9647fbb06cd6fc529deec6f3dda82}.
\item The supplied attribution identifies V. S. Monakhov as proposer, with proposed year 2002, and cites \emph{The Kourovka Notebook (Unsolved Problems in Group Theory)}, 15th issue (2002), Problem 15.61; it also identifies the 21st issue edition (2026), arXiv:1401.0300v48. The collection description names E. I. Khukhro and V. D. Mazurov as editors.
\item The registry supplies the status label \texttt{open} and reports the affirmative special case concerning proper supersoluble subgroups. These are supplied assessments, not independently verified mathematical or literature conclusions. No separate citation for the special case or literature-verification date is supplied.
\end{itemize}
\end{document}
